\textbf{1. Kiến thức cần nhớ}
\begin{itemize}
\item \textbf{Đa giác đều} là đa giác có tất cả các cạnh bằng nhau và tất cả các góc bằng nhau.
\begin{itemize}
\item Tam giác đều ($n = 3$), hình vuông ($n = 4$), ngũ giác đều ($n = 5$), lục giác đều ($n = 6$)\dots
\item Số đo mỗi góc của đa giác đều $n$ cạnh: $\boxed{\alpha = \dfrac{(n - 2) \cdot 180^\circ}{n}}$.
\end{itemize}
\item Bất kì đa giác đều nào cũng có một và chỉ một đường tròn ngoại tiếp, một và chỉ một đường tròn nội tiếp và hai đường tròn này \textbf{đồng tâm} (gọi là tâm của đa giác đều).
\end{itemize}

\textbf{2. Phương pháp giải các dạng}

%% DANG 9H9?30-1
\textbf{Dạng 1. Nhận biết đa giác đều, tính góc}
\begin{enumerate}
\item Áp dụng công thức tính số đo góc: Tam giác đều: $60^\circ$; hình vuông: $90^\circ$; ngũ giác đều: $108^\circ$; lục giác đều: $120^\circ$.
\item Góc ở tâm của đa giác đều $n$ cạnh: $\boxed{\beta = \dfrac{360^\circ}{n}}$.
\end{enumerate}

%% DANG 9H9?30-2
\textbf{Dạng 2. Phép quay giữ nguyên hình đa giác đều}
\begin{enumerate}
\item Phép quay tâm $O$ (tâm đa giác đều) góc quay $\alpha = k \cdot \dfrac{360^\circ}{n}$ ($k \in \mathbb{Z}$) biến đa giác đều thành chính nó.
\item Lục giác đều: góc quay $60^\circ, 120^\circ, 180^\circ\dots$ giữ nguyên hình.
\end{enumerate}

%% DANG 9H9?30-3
\textbf{Dạng 3. Đa giác đều trong thực tế}
\begin{enumerate}
\item Ứng dụng trong kiến trúc, tổ ong (lục giác đều tiết kiệm vật liệu nhất), hoa văn gạch lát nền (chỉ có tam giác đều, hình vuông, lục giác đều lát kín mặt phẳng mà không để lại khe hở vì các góc là ước của $360^\circ$).
\end{enumerate}
